% ===========================================================================
\section{Lower bounds: is lossless even possible?}\label{sec:lower}
% ===========================================================================

\paragraph{Framing.}
For a \emph{fixed} teacher, whether lossless distillation is possible is purely a question of
approximation theory: is $p$ in the closure of $\{q_\theta\}$? Information theory enters once we ask
what a class of $2^R$ students can do \emph{uniformly over teachers}, either on average under a prior
$\pi$ over teachers or in the worst case over a class. We then view the student as a code of $R$ bits for
the teacher's conditional law, with distortion $\E_x\KL(p(\cdot\mid x)\|q_\theta(\cdot\mid x))$. The random
teacher is denoted $\mathbf T\sim\pi$, and its conditionals $p_{\mathbf T}(\cdot\mid x)$.

\subsection{An information-theoretic capacity floor}

\begin{theorem}[Capacity floor]\label{thm:capacity}
\Pa{\cref{as:finite}}\Ck Let $\theta^\star(\mathbf T)$ be any measurable choice of student from a class
of at most $2^R$ elements. Fix a design $x_{1:N}\in\cX^N$ and a replication number $m\ge1$. Let
$Y=(Y_{ij})_{i\le N,j\le m}$ with $Y_{ij}\sim p_{\mathbf T}(\cdot\mid x_i)$ independent given $\mathbf T$. Then
\[
\E_\pi\,\frac1N\sum_{i=1}^N\KL\big(p_{\mathbf T}(\cdot\mid x_i)\,\|\,q_{\theta^\star(\mathbf T)}(\cdot\mid x_i)\big)\;\ge\;\frac{I(\mathbf T;Y)-R\ln2}{mN}.
\]
In particular:
(a)~if $\cD$ is uniform on $\{x_1,\dots,x_N\}$, the left side is $\E_\pi\E_{x\sim\cD}\KL$;
(b)~for general $\cD$,
$\E_\pi\E_{x\sim\cD}\KL\ge\sup_{N,m}\frac{I(\mathbf T;Y\mid X_{1:N})-R\ln2}{mN}$ with $X_i\overset{\text{iid}}\sim\cD$.
\end{theorem}
\begin{proof}
Write $\Theta=\theta^\star(\mathbf T)$. KL tensorises, so
$m\sum_i\KL(p_{\mathbf T}(\cdot\mid x_i)\|q_\Theta(\cdot\mid x_i))=\KL(\otimes_{ij}p_{\mathbf T}(\cdot\mid x_i)\,\|\,\otimes_{ij}q_\Theta(\cdot\mid x_i))
=\E[-\log q_\Theta(Y)\mid\mathbf T]-H(Y\mid\mathbf T)$.
Average over $\mathbf T$. Since $q_\Theta$ is a valid conditional law of $Y$ given $\Theta$, the
cross-entropy satisfies $\E[-\log q_\Theta(Y)]\ge H(Y\mid\Theta)$. Hence
\[
m\,\E\sum_i\KL_i\ge H(Y\mid\Theta)-H(Y\mid\mathbf T)=I(Y;\mathbf T)-I(Y;\Theta)\ge I(Y;\mathbf T)-H(\Theta)\ge I(Y;\mathbf T)-R\ln2,
\]
where the equality uses $H(Y)$ on both sides and the fact that $\Theta$ is a function of $\mathbf T$.
For (b), apply the bound conditionally on $X_{1:N}$ and average, since
$\E_{X_{1:N}}\frac1N\sum_i\KL(\cdot;X_i)=\E_{x\sim\cD}\KL$.
\end{proof}

\noindent
Let $I_\infty:=\sup_{N,m}I(\mathbf T;Y\mid X_{1:N})$ be the \emph{functional information} of the teacher
on $\cD$: how much its input--output behaviour reveals about which teacher it is. \Cref{thm:capacity}
says that on average over $\pi$, lossless distillation is \emph{impossible} whenever $R\ln2<I_\infty$. The
student's description length must be at least the teacher's functional information. This is not the
teacher's parameter count: trained networks are highly redundant as functions, which is what makes
pruning, quantisation and distillation work at all.

\begin{corollary}[$K$ independent facts; the floor is tight]\label{cor:facts}
\Pv\Ck Let the teacher be a uniformly random function $F:[K]\to\cV$ with $p(\cdot\mid i)=\delta_{F(i)}$, and let
$\cD$ be uniform on $[K]$. Every $R$-bit student class satisfies
$\E\,\E_{x\sim\cD}\KL\ge\big(\ln V-\frac{R\ln2}K\big)_+$. The student that memorises
$\lfloor R/\log_2V\rfloor$ facts and is uniform elsewhere attains this up to rounding.
\end{corollary}
\begin{proof}
Take $N=K$, $x_i=i$, $m=1$: $I(F;Y)=K\ln V$. The memorising student has KL $\ln V$ on each fact it does not
store and $0$ on the others.
\end{proof}

\begin{corollary}[Behavioural floor for $K$ facts]\label{cor:facts-tv}
\Pv In the setting of \cref{cor:facts} with $V\ge2$, let $\bar e:=\E\,\E_{x\sim\cD}\TV\big(p(\cdot\mid x),q_{\theta^\star}(\cdot\mid x)\big)$.
Every $R$-bit student class satisfies
\[
h(\bar e)+\bar e\ln(V-1)\ \ge\ \ln V-\frac{R\ln2}K .
\]
The left side increases from $0$ to $\ln V$ as $\bar e$ goes from $0$ to $1-\frac1V$. So $\bar e>0$ whenever $R<K\log_2V$, and
$\bar e\ge1-\frac1V$ at $R=0$, which is tight. In particular $\bar e\ge(\ln V-\ln2-R\ln2/K)/\ln(V-1)$ for $V\ge3$. The KL
floor of \cref{cor:facts} is therefore not an artefact of KL: behavioural losslessness also fails once
$R<K\log_2V$.
\end{corollary}
\begin{proof}
Let $\hat G_i\sim q_{\theta^\star}(\cdot\mid i)$ independently given $\theta^\star$. Then
$e_i:=\P(\hat G_i\ne F(i))=\E\TV_i$, because the teacher is a point mass. Since the $F(i)$ are independent and
$\hat G$ depends on $F$ only through $\theta^\star$,
$\sum_iI(F(i);\hat G_i)\le I(F;\hat G)\le H(\theta^\star)\le R\ln2$. Fano gives
$I(F(i);\hat G_i)\ge\ln V-h(e_i)-e_i\ln(V-1)$. Sum over $i$ and use concavity of $h$:
$K[\ln V-h(\bar e)-\bar e\ln(V-1)]\le R\ln2$. Monotonicity follows from
$\frac{\mathrm d}{\mathrm de}[h(e)+e\ln(V-1)]=\ln\frac{(1-e)(V-1)}e\ge0$ for $e\le1-\frac1V$. The explicit bound uses $h\le\ln2$.
\end{proof}

\begin{corollary}[Metric-entropy version]\label{cor:metric}
\Pv Suppose $\Theta_\TS$ has an $\eps$-cover $\Theta_\eps$ in $\|\log q_\theta-\log q_{\theta'}\|_\infty$ (over the
design and $\cV^T$). Then $\KL(p\|q_{\theta'})\le\KL(p\|q_\theta)+\eps$ whenever $\theta'$ is $\eps$-close to
$\theta$. Hence
\[
\E_\pi\min_{\theta\in\Theta_\TS}\frac1N\sum_i\KL_i\ge\frac{I(\mathbf T;Y)-\log\cN(\eps)}{mN}-\eps.
\]
For $P$ parameters in $[-B_\theta,B_\theta]^P$ with $\theta\mapsto\log q_\theta$ being $L$-Lipschitz in $\ell_\infty$,
$\log\cN(\eps)\le P\log\lceil LB_\theta/\eps\rceil$, so the effective rate is
$R_{\mathrm{eff}}(\eps)\approx P\log_2(LB_\theta/\eps)$ bits. Pseudo-dimension and Rademacher bounds enter
through the usual covering-number estimates.
\end{corollary}
\begin{proof}
$\KL(p\|q_{\theta'})-\KL(p\|q_\theta)=\E_p[\log q_\theta-\log q_{\theta'}]\le\eps$. Apply \cref{thm:capacity} to
the finite class $\Theta_\eps$.
\end{proof}

\begin{remark}[Bits per parameter]
The rate that matters is \emph{usable} bits. \citet{allenzhu2024physics} report, in controlled synthetic
settings, a knowledge capacity of about $2$ bits per parameter for language models. \citet{kumar2024scaling}
model how precision changes the effective parameter count. Combining such estimates with
\cref{thm:capacity} gives a heuristic necessary condition. A student with $N_\TS$ parameters cannot be
lossless on content whose functional information exceeds roughly $2N_\TS$ bits. Treat this as an
order-of-magnitude guide; the constant is an empirical estimate for a particular kind of knowledge.
\end{remark}

\subsection{Distillation as rate--distortion}

\begin{definition}[Distillation distortion--rate function]
For a prior $\pi$ on teachers, a reproduction alphabet of student functions, and distortion
$\mathsf d(\mathbf T,\hat{\mathbf T})=\E_{x\sim\cD}\KL(p_{\mathbf T}(\cdot\mid x)\|p_{\hat{\mathbf T}}(\cdot\mid x))$, define
$D_\pi(R):=\inf\{\E\,\mathsf d(\mathbf T,\hat{\mathbf T}):\ I(\mathbf T;\hat{\mathbf T})\le R\ln2\}$.
\end{definition}

\begin{theorem}[One-shot converse]\label{thm:rd}
\Pv Every class of $2^R$ students satisfies $\E_\pi\min_\theta\mathsf d(\mathbf T,q_\theta)\ge D_\pi(R)$.
\end{theorem}
\begin{proof}
$\theta^\star(\mathbf T)$ is a deterministic function of $\mathbf T$ taking at most $2^R$ values, so
$I(\mathbf T;\theta^\star)\le H(\theta^\star)\le R\ln2$. It is therefore feasible in the infimum.
\end{proof}
\noindent
Unlike Shannon's asymptotic theorem \citep{shannon1959coding}, no block coding is needed. The converse
holds for a single teacher draw because it uses only $I\le H\le\log(\#\text{codewords})$
\citep{cover2006elements,polyanskiy2024information}. For log-loss distortion the rate--distortion
function has the simple form $R(D)=H-D$ \citep{courtade2014multiterminal}. \Cref{cor:facts} is an
instance of this.

\begin{proposition}[Power-law spectrum gives a power-law floor]\label{prop:powerlaw}
\Mo\Ck \emph{Model.} Near-lossless regime (\cref{prop:local}), in Fisher-whitened centred
log-probability coordinates, where $\KL\approx\frac12\|z-\hat z\|^2$. The teacher's whitened
log-probability function is a centred Gaussian vector with covariance eigenvalues $\lambda_1\ge\lambda_2\ge\cdots$.
Reverse water-filling \citep[Thm.~10.3.3]{cover2006elements} gives
\[
D(R)=\tfrac12\sum_i\min(w,\lambda_i),\qquad\sum_i\tfrac12\log_+\frac{\lambda_i}w=R\ln2 .
\]
If $\lambda_i=c\,i^{-\alpha}$ with $\alpha>1$, then
$D(R)=\frac c2\,\frac\alpha{\alpha-1}\big(\tfrac{2R\ln2}\alpha\big)^{-(\alpha-1)}(1+o(1))$, so $D(R)\asymp R^{-(\alpha-1)}$.
\end{proposition}
\begin{proof}[Proof sketch]
With $k$ active components, $w=\lambda_k$. The rate is
$\frac\alpha2\sum_{i\le k}\log\frac ki=\frac{\alpha}2(k-O(\log k))$ by Stirling, so $k\approx2R\ln2/\alpha$. The
distortion is $\frac12[kw+\sum_{i>k}\lambda_i]=\frac c2k^{1-\alpha}\big(1+\frac1{\alpha-1}\big)(1+o(1))$.
The fitted exponents in \texttt{check\_capacity.py} match $-(\alpha-1)$ to within $0.04$ for
$\alpha\in\{1.5,2,3\}$.
\end{proof}

\begin{proposition}[Finite teacher dimension: from power law to exponential]\label{prop:phase}
\Mo\Ck If $\lambda_i=0$ for $i>k_\TT$, then as soon as $w<\lambda_{k_\TT}$,
$D(R)=\frac{k_\TT}2\big(\prod_{i\le k_\TT}\lambda_i\big)^{1/k_\TT}e^{-2R\ln2/k_\TT}$. The distortion decays
\emph{exponentially} in the student's rate once the student can resolve every component the teacher has.
For $k_\TT=1000$ and $\alpha=2$, $D$ drops by a factor $e^{4}\approx55$ when $R\ln2$ doubles from $2000$ to
$4000$ nats. Below $R\ln2\approx k_\TT$ nats it follows the power law.
\end{proposition}

\begin{proposition}[No capacity gap under optimal rate-limited allocation]\label{prop:nogap}
\Mo Let the ground truth have spectrum $(\lambda_i)_{i\ge1}$ and let a teacher of size $k_\TT$ know the first
$k_\TT$ components exactly. The best $R$-bit student's distortion \emph{with respect to the ground truth} is
non-increasing in $k_\TT$.
\end{proposition}
\begin{proof}
Let $E(k_\TT)=D_{k_\TT}(R)+\frac12\sum_{i>k_\TT}\lambda_i$. For $k'>k_\TT$, one feasible code for the larger
teacher uses the optimal code for the first $k_\TT$ components and reproduces $0$ elsewhere. Hence
$D_{k'}(R)\le D_{k_\TT}(R)+\frac12\sum_{k_\TT<i\le k'}\lambda_i$, so $E(k')\le E(k_\TT)$.
\end{proof}

\subsection{What this implies for distillation scaling laws}

Supervised pre-training follows power laws in parameters $N$ and tokens $D$ \citep{kaplan2020scaling}. A
widely used joint form is the additive law $L(N,D)=E+A/N^{a}+B/D^{b}$ of \citet{hoffmann2022training}. \citet{busbridge2025distillation} fit a distillation
scaling law in which the student's cross-entropy depends on the teacher's cross-entropy and on the
student's supervised-scaling loss. That law captures a ``capacity gap'': beyond some point, stronger
teachers yield worse students \citep[see also][]{cho2019efficacy,mirzadeh2020improved,zhang2023capacity}.
We are confident of this qualitative content but not of the exact parametrisation in
\citet{busbridge2025distillation}. The claim of \citet{zhang2023capacity} that the optimal teacher is
about $2.5\times$ the student is quoted from memory and should be checked. The results above imply the
following.

\begin{heuristic}[Student-size exponent]\label{heur:size}
With a constant number of usable bits per parameter, $R_\TS\propto N_\TS$. The approximation part of the
gap should then scale as $N_\TS^{-(\alpha-1)}$, where $\alpha$ is the decay exponent of the teacher's
\emph{Fisher-weighted} centred log-probability spectrum. This spectrum is measurable (\cref{sec:experiments}).
It should cross over to faster (exponential) decay once $R_\TS\ln2$ approaches the teacher's functional
information (\cref{prop:phase}).
\end{heuristic}

\begin{heuristic}[Data exponent: soft labels are steeper]\label{heur:data}
Consider a projection-estimator model. The target has coefficients of energy $\lambda_i=ci^{-\alpha}$ in
an orthonormal basis under $\cD$, and the student estimates the first $k$ from $n$ samples. The error is
approximately $\sigma^2\frac kn+\sum_{i>k}\lambda_i$, where $\sigma^2$ is the effective noise.
With hard labels, $\sigma^2$ includes label noise $\sigma^2_0>0$. The optimum is $k\asymp n^{1/\alpha}$, giving
error $\asymp n^{-(\alpha-1)/\alpha}$.
With soft labels, the only noise is the aliased tail, $\sigma^2\asymp k^{1-\alpha}$. Then $k\asymp n$ and the
error is $\asymp n^{-(\alpha-1)}$.
The rigorous version of this noiseless/noisy crossover in kernel regression is due to
\citet{cui2021generalization}. The prediction is that the distillation-token exponent should be
steeper for full-logit distillation than for sequence-level (hard-label) distillation, with top-$k$
logits in between.
\end{heuristic}

\begin{heuristic}[Capacity gap as a finite-sample aliasing effect]\label{heur:gap}
In the same model, let the teacher have $k_\TT$ components and the student $k<k_\TT$. With soft labels,
the student's error with respect to the ground truth is approximately
\[
\frac kn\sum_{k<i\le k_\TT}\lambda_i+\sum_{i>k}\lambda_i .
\]
The first term increases in $k_\TT$ (teacher content the student cannot represent aliases into the
components it estimates). For $k_\TT\le k$ the error is $\sum_{i>k_\TT}\lambda_i$, which decreases in $k_\TT$.
Hence the optimal teacher has $k_\TT\approx k$, and the penalty for a stronger teacher is at most a factor
$(1+k/n)$ on the tail. It vanishes as $n\to\infty$. Together with \cref{prop:nogap}, this predicts that
capacity gaps are \emph{finite-data and optimisation effects}, and should shrink as distillation tokens
grow at fixed student size.
\end{heuristic}

\begin{conjecture}[Distillation gap law]\label{conj:law}
For a teacher whose Fisher-weighted log-probability spectrum decays as $i^{-\alpha}$, the per-token
forward-KL gap of a well-optimised student with $R_\TS$ usable bits and $n$ distillation tokens satisfies
\[
\eoff^{\KL}\asymp A\,R_\TS^{-(\alpha-1)}+B\,n^{-b}+\frac{k(R_\TS)}n\,\mathrm{Tail}_\TT\big(k(R_\TS)\big),
\]
with $b=\alpha-1$ for full soft labels and $b=(\alpha-1)/\alpha$ for hard labels. The first term switches
to exponential decay once $R_\TS\ln2\gtrsim I_\infty$.
\end{conjecture}
